\medskip

\emph{Pour chacune des affirmations suivantes, indiquer si elle est vraie ou fausse. Chaque réponse doit être justifiée. Une réponse non justifiée ne rapporte aucun point.}

\medskip

Dans l'espace muni d'un repère orthonormé, on considère les points suivants :

\begin{center}
	A$(2~;~0 ~;~0)$,\qquad B$(0 ~;~4~;~3)$,\qquad C$(4 ~;~4 ~;~1)$,\qquad D$(0 ~;~0 ~;~4)$\quad et \quad H$(-1~;~1~;~2)$
\end{center}

\smallskip

\textbf{Affirmation 1 :} les points A, C et D définissent un plan $\mathcal{P}$ d'équation $8 x- 5 y+ 4z - 16 = 0$.

\medskip

\textbf{Affirmation 2 :} les points A, B, C et D sont coplanaires.

\medskip

\textbf{Affirmation 3 :} les droites (AC) et (BH) sont sécantes.

\bigskip

On admet que le plan (ABC) a pour équation cartésienne $x - y + 2z - 2 = 0$.

\smallskip

\textbf{Affirmation 4 :} le point H est le projeté orthogonal du point D sur le plan (ABC).
